\section*{Glosario y abreviaturas}
\addcontentsline{toc}{section}{Glosario y abreviaturas}
\markboth{\MakeUppercase{Glosario y abreviaturas}}{}

\begin{description}[style=nextline, leftmargin=1.6cm]
  \item[BA] Bases Administrativas del proceso TFEP-01/2026.
  \item[BTT] Bases Técnicas Transversales TFEP-01/2026.
  \item[CP] Bases Técnicas del Caso (Caso 06 --- Portuaria).
  \item[RT] Código de requisito de las Bases Técnicas Transversales
    (por ejemplo, \texttt{RT-26.01}).
  \item[RF] Código de requerimiento funcional del catálogo del
    Subdocumento 3 (por ejemplo, \texttt{RF-REF-08}).
  \item[Alcance 1, 2 y 3] Categorías de emisiones de gases de efecto
    invernadero: directas de la organización, indirectas por energía
    adquirida e indirectas de la cadena de valor, respectivamente.
  \item[CO\textsubscript{2}e] Dióxido de carbono equivalente. Unidad que
    expresa el efecto de distintos gases de efecto invernadero en una
    escala común.
  \item[EDIFACT] Estándar internacional de mensajería electrónica para el
    intercambio de datos entre navieras y terminales.
  \item[GLEC Framework] Marco del \emph{Global Logistics Emissions Council}
    para la contabilización y el reporte de emisiones logísticas.
  \item[Intensidad de carbono] Emisión atribuida por unidad de actividad.
    En este subdocumento, kilogramos de CO\textsubscript{2}e por contenedor
    movilizado.
  \item[Marcha blanca] Período de operación supervisada con datos y usuarios
    reales, en convivencia con la operación vigente y con plan de reversión
    activo, previo al paso a producción.
  \item[Ralentí] Tiempo en que un equipo mantiene el motor encendido sin
    ejecutar trabajo productivo.
  \item[\emph{Reefer}] Contenedor refrigerado. Por extensión, el patio y las
    tomas de conexión que lo alimentan.
  \item[Remoción] Movimiento de patio que reubica un contenedor sin
    trasladarlo a su destino, ejecutado para acceder a otro que está debajo.
  \item[Sello de tiempo] Constancia criptográfica de que un dato existía en
    un instante determinado, verificable por un tercero.
  \item[Simulación de eventos discretos] Técnica de modelado que representa
    la operación como una secuencia de eventos en el tiempo, para evaluar
    escenarios sin intervenir el sistema real.
  \item[Toma] Punto de conexión eléctrica del patio refrigerado al que se
    enchufa un contenedor.
  \item[TRL] \emph{Technology Readiness Level}, nivel de madurez tecnológica.
    En este documento se declara según la escala de la norma ISO 16290:2013.
  \item[Ventana de atraque] Franja horaria comprometida con la naviera para
    la atención de una nave.
\end{description}

\vspace{4pt}
\noindent\small\textcolor{terabyteGris}{%
Agregar cada sigla o término técnico nuevo apenas se use por primera vez
en el cuerpo del documento.}
